\chapter{\ifenglish Background Knowledge and Theory\else ทฤษฎีที่เกี่ยวข้อง\fi}

เนื้อหาในบทนี้อธิบายแนวคิดและเทคโนโลยีที่เกี่ยวข้องกับโครงงาน
ได้แก่ ระบบ SIEM แหล่งข้อมูล Logs มาตรฐานการจัดรูปแบบข้อมูล กรอบความรู้ MITRE ATT\&CK
โมเดลภาษาขนาดใหญ่ เทคนิค RAG และตัวชี้วัดที่ใช้ประเมินผล
เพื่อเป็นพื้นฐานสำหรับการทำความเข้าใจการออกแบบในบทถัดไป

\section{Security Information and Event Management (SIEM)}

SIEM คือระบบที่รวบรวม Logs และเหตุการณ์ด้านความปลอดภัย (Security Events)
จากแหล่งต่าง ๆ ภายในองค์กร เช่น เครื่องลูกข่าย เซิร์ฟเวอร์ ไฟร์วอลล์ และแอปพลิเคชัน
มาจัดเก็บไว้ในที่เดียว เพื่อให้สามารถค้นหา เชื่อมโยง (Correlate) และสร้างการแจ้งเตือน (Alert)
เมื่อพบพฤติกรรมที่ตรงกับกฎการตรวจจับ (Detection Rule)
SIEM จึงเป็นเครื่องมือหลักของศูนย์ปฏิบัติการด้านความปลอดภัย (Security Operations Center: SOC)

ในโครงงานนี้ใช้ \textbf{Splunk Enterprise} เป็นระบบ SIEM
โดย Logs ถูกจัดเก็บเป็นดัชนี (Index) และสามารถเรียกค้นได้ด้วยภาษา Search Processing Language (SPL)
หรือผ่าน REST API ซึ่ง Pipeline ของโครงงานใช้สำหรับดึง Logs ที่เกี่ยวข้องกับ Alert

\subsection{SOC Analyst}
SOC Analyst คือเจ้าหน้าที่ด้านความปลอดภัยไซเบอร์ที่ทำหน้าที่เฝ้าระวัง ตรวจสอบ คัดกรอง
และวิเคราะห์เหตุการณ์ด้านความปลอดภัย เพื่อระบุภัยคุกคามและสนับสนุนการตอบสนองต่อเหตุการณ์ (Incident Response)
ซึ่งเป็นผู้มีส่วนได้ส่วนเสีย (Stakeholder) หลักของโครงงานนี้

\section{แหล่งข้อมูล Security Logs}

\subsection{Windows Sysmon}
System Monitor (Sysmon)~\cite{sysmon} เป็นบริการของ Microsoft Sysinternals
ที่บันทึกกิจกรรมของระบบปฏิบัติการ Windows อย่างละเอียดลงใน Windows Event Log
โดยแต่ละประเภทกิจกรรมมีหมายเลข Event ID เฉพาะ ตัวอย่าง Event ID ที่สำคัญต่อโครงงานแสดงในตารางที่~\ref{tab:sysmon-eid}
Events ที่เกี่ยวข้องกับ Process จะมีตัวระบุ \texttt{ProcessGuid} ที่ไม่ซ้ำกัน
และ Event ID 1 จะบันทึก \texttt{ParentProcessGuid} ไว้ด้วย
จึงสามารถนำมาสร้างความสัมพันธ์แบบ Parent--Child ระหว่าง Process ได้

\begin{table}[ht]
  \centering
  \caption{ตัวอย่าง Sysmon Event ID ที่ใช้ในโครงงาน}
  \label{tab:sysmon-eid}
  \begin{tabular}{|c|l|l|}
    \hline
    \textbf{Event ID} & \textbf{ชื่อเหตุการณ์} & \textbf{ข้อมูลสำคัญ} \\ \hline
    1  & Process Creation        & Image, CommandLine, ParentImage, User \\ \hline
    3  & Network Connection      & SourceIp, DestinationIp, DestinationPort \\ \hline
    10 & Process Access          & SourceImage, TargetImage, GrantedAccess \\ \hline
    11 & File Create             & Image, TargetFilename \\ \hline
    23 & File Delete (archived)  & Image, TargetFilename \\ \hline
  \end{tabular}
\end{table}

\subsection{Linux Auditd}
Linux Audit Daemon (auditd)~\cite{auditd} เป็นระบบบันทึกเหตุการณ์ของ Linux Kernel
เช่น การเรียกใช้ System Call การรันคำสั่ง และการเข้าถึงไฟล์
โดยหนึ่งเหตุการณ์อาจประกอบด้วยหลาย Record (เช่น \texttt{SYSCALL}, \texttt{EXECVE}, \texttt{PATH})
ที่เชื่อมกันด้วยหมายเลข Event เดียวกัน
ในโครงงานนี้ Auditd อยู่ในขอบเขตการออกแบบ ส่วนการทดลองใช้ Windows Sysmon

\section{Common Information Model (CIM)}
เนื่องจาก Logs จากแต่ละแหล่งมีรูปแบบและชื่อ Field แตกต่างกัน
(เช่น Sysmon เป็น XML ส่วน Auditd เป็นข้อความแบบ Key=Value)
Splunk จึงกำหนด Common Information Model (CIM)~\cite{splunk-cim}
ซึ่งเป็นมาตรฐานชื่อ Field และโครงสร้างข้อมูลร่วม เช่น \texttt{process}, \texttt{parent\_process}, \texttt{user}, \texttt{dest}
% การ Normalize Logs ให้อยู่ในรูปแบบร่วมทำให้ขั้นตอนถัดไปของ Pipeline ทำงานได้กับ Logs หลายแหล่งโดยไม่ต้องเขียนตรรกะแยก
โครงงานนำแนวคิดของ CIM มาใช้ในการกำหนดโครงสร้าง Field กลางของโครงงาน เพื่อให้ขั้นตอนถัดไปของ Pipeline ทำงานกับ Logs หลายแหล่งได้ด้วยตรรกะเดียวกัน

\section{MITRE ATT\&CK}
MITRE ATT\&CK~\cite{mitre-attack,strom2018attack} เป็นฐานความรู้ระดับสากลที่รวบรวมพฤติกรรมของผู้โจมตีไซเบอร์จากเหตุการณ์จริง
โดยจำแนกเป็น \textbf{Tactic} (เป้าหมายของการโจมตี เช่น Discovery, Credential Access)
\textbf{Technique} (วิธีการที่ใช้เพื่อบรรลุเป้าหมาย) และ \textbf{Sub-technique} (รายละเอียดเฉพาะของวิธีการ)
แต่ละเทคนิคมีรหัสประจำ เช่น T1087 Account Discovery และ T1087.001 Account Discovery: Local Account~\cite{mitre-t1087-001}
ซึ่งผู้โจมตีพยายามดึงรายชื่อบัญชีผู้ใช้ภายในเครื่อง เช่น ใช้คำสั่ง \texttt{net user} หรือ \texttt{net localgroup}

ATT\&CK Enterprise เวอร์ชัน 19.2 ที่ใช้ในโครงงานประกอบด้วย 15 Tactics, 222 Techniques
และ 475 Sub-techniques (ไม่นับรายการที่ถูกยกเลิก) และมีการปรับปรุงเวอร์ชันประมาณปีละ 2 ครั้ง
ในโครงงานนี้ ATT\&CK ถูกใช้ใน 2 บทบาท คือ (1) เป็นป้ายกำกับ (Label) ของ Ground Truth สำหรับประเมินผล
และ (2) เป็นฐานความรู้ (Knowledge Base) ให้ LLM อ้างอิงผ่าน RAG

\section{ชุดข้อมูลและแนวคิด Scenario Dataset}

\subsection{OTRF Security-Datasets}
OTRF Security-Datasets~\cite{otrf-security-datasets} เป็นชุดข้อมูลสาธารณะที่บันทึก Logs จากการจำลองการโจมตี
โดยมีคำอธิบายมุมมองผู้โจมตี (Adversary View) ระบุคำสั่งและการกระทำของผู้โจมตีอย่างละเอียด
% จึงเหมาะสำหรับใช้เป็น Benchmark ในการออกแบบและวัดผล Pipeline
จึงใช้เป็นชุดข้อมูลสำหรับทดสอบ (Test Set) ซึ่งประเมินผลเพียงครั้งเดียวหลังกำหนดค่าพารามิเตอร์แล้ว

\subsection{Splunk Attack Data}
Splunk Attack Data~\cite{splunk-attack-data} เป็นคลังชุดข้อมูลที่จัดตามรหัสเทคนิค ATT\&CK
แต่ละไฟล์ระบุว่าเป็นเทคนิคใด แต่ไม่ได้ระบุว่า Event ใดเป็นส่วนของการโจมตี
% โครงงานจึงใช้ชุดข้อมูลนี้เป็น Validation โดย Replay Logs เข้า Splunk Enterprise เพื่อทดสอบในสภาพใกล้เคียงการใช้งานจริง
โครงงานจึงใช้ชุดข้อมูลนี้เป็นชุดข้อมูลสำหรับพัฒนา (Development Set) โดยระบุ Event สำคัญเพิ่มเติม (Event-level Annotation) เพื่อใช้ปรับพารามิเตอร์ของ Pipeline และ Replay Logs เข้า Splunk Enterprise เพื่อทดสอบในสภาพใกล้เคียงการใช้งานจริง

\subsection{นิยามที่ใช้ในโครงงาน}
\begin{itemize}
  \item \textbf{Attack Scenario} -- การจำลองการโจมตีหนึ่งสถานการณ์ในสภาพแวดล้อมการทดลอง
  \item \textbf{Scenario Dataset} -- ชุดของ Events ที่เกิดขึ้นในบริบทของ Attack Scenario หนึ่ง
  \item \textbf{Event / Log Event} -- บันทึกหนึ่ง Record ของกิจกรรม เช่น Process Creation หรือ Network Connection
  \item \textbf{Ground Truth} -- เทคนิค ATT\&CK ที่ใช้สร้าง Scenario นั้น
        เป็น Ground Truth ระดับ Scenario/Dataset ไม่ใช่ Label ของทุก Event
        (1 Scenario Dataset $\rightarrow$ 1 Ground-truth Technique $\rightarrow$ หลาย Events)
  \item \textbf{Noise} -- เหตุการณ์ใน Logs ที่ไม่เกี่ยวข้องกับการวิเคราะห์ รวมถึง Field และ Event ที่ไม่จำเป็น
  \item \textbf{Investigation Context} -- ชุดเหตุการณ์และความรู้ที่เกี่ยวข้องกับเหตุการณ์ที่น่าสงสัย
        ซึ่งถูกจัดเตรียมเพื่อให้ LLM ใช้เหตุผลกับเหตุการณ์นั้น
\end{itemize}

\section{Large Language Model (LLM)}
LLM คือโมเดลภาษาที่สร้างจากสถาปัตยกรรม Transformer~\cite{vaswani2017attention}
และฝึกด้วยข้อความจำนวนมหาศาล ทำให้สามารถเข้าใจและสร้างข้อความ รวมถึงใช้เหตุผลกับข้อมูลที่ได้รับได้
LLM ประมวลผลข้อความในหน่วย \textbf{Token} และรับข้อมูลได้จำกัดตามขนาด \textbf{Context Window}
ต้นทุนและเวลาในการประมวลผลจึงแปรผันตามจำนวน Token ที่ส่งเข้าไป
นอกจากนี้ งานวิจัยพบว่าเมื่อ Context ยาวขึ้น LLM มีแนวโน้มใช้ข้อมูลที่อยู่ช่วงกลางของ Context ได้ไม่ดี~\cite{liu2024lost}
การคัดเลือกและจัดระเบียบข้อมูลก่อนส่งให้ LLM จึงช่วยลด Noise และทำให้กลไก Attention ของโมเดลโฟกัสข้อมูลที่สำคัญได้ดีขึ้น

\subsection{Local LLM: Qwen3.5 และ Ollama}
โครงงานเลือกใช้ LLM ขนาดเล็กที่ทำงานภายในเครื่อง ได้แก่ Qwen3.5 ขนาด 9B พารามิเตอร์
ซึ่งเป็นโมเดลในตระกูล Qwen~\cite{qwen3} ที่เผยแพร่แบบ Open-weight
และรันผ่าน Ollama~\cite{ollama} ที่ให้บริการโมเดลผ่าน API ภายในเครื่อง
ข้อดีของ Local LLM คือ Logs ขององค์กรไม่ต้องถูกส่งออกไปภายนอก และไม่มีค่าบริการ API ตามจำนวน Token
แต่โดยทั่วไปมีความสามารถจำกัดกว่าโมเดลเชิงพาณิชย์ขนาดใหญ่
เป้าหมายของโครงงานจึงเป็นการแสดงว่า Context ที่ดีช่วยให้โมเดลขนาดเล็กตอบได้ถูกต้อง

\section{Retrieval-Augmented Generation (RAG)}
RAG~\cite{lewis2020rag} คือเทคนิคที่เสริมความรู้ให้ LLM ขณะตอบคำถาม
โดยแบ่งเอกสารความรู้เป็นส่วนย่อย (Chunk) แล้วแปลงแต่ละ Chunk เป็นเวกเตอร์ด้วย Embedding Model
เก็บไว้ในฐานข้อมูลเวกเตอร์ (Vector Database)
เมื่อมีคำถาม (Query) ระบบจะแปลง Query เป็นเวกเตอร์และค้นหา Chunk ที่มีความคล้ายสูงสุด $k$ อันดับ (Top-$k$)
ด้วยการวัดความคล้ายแบบ Cosine Similarity ดังสมการ~\eqref{eq:cosine}
แล้วนำ Chunk เหล่านั้นใส่ใน Prompt เป็น Context ให้ LLM

\begin{equation}\label{eq:cosine}
\mathrm{sim}(\mathbf{q},\mathbf{d}) = \frac{\mathbf{q}\cdot\mathbf{d}}{\lVert\mathbf{q}\rVert\,\lVert\mathbf{d}\rVert}
\end{equation}

\noindent
โดยที่ $\mathbf{q}$ คือเวกเตอร์ของ Query และ $\mathbf{d}$ คือเวกเตอร์ของ Chunk ในฐานความรู้

\section{รูปแบบการนำเสนอข้อมูล (Data Representation)}
ข้อมูลชุดเดียวกันสามารถเขียนได้หลายรูปแบบ ซึ่งส่งผลต่อจำนวน Token และความเข้าใจของ LLM
รูปแบบที่โครงงานพิจารณา ได้แก่
\textbf{XML} (รูปแบบต้นฉบับของ Sysmon มีแท็กซ้ำซ้อนจำนวนมาก),
\textbf{JSON} (โครงสร้างชัดเจนและเป็น Output ของขั้นตอน Normalize),
\textbf{TOON} (Token-Oriented Object Notation~\cite{toon-format} ที่ตัดเครื่องหมาย \{\}, "" ออกเพื่อลด Token),
\textbf{Key=Value} (รูปแบบที่ Splunk ใช้แสดงผล) และ
\textbf{Timeline} (สรุปเหตุการณ์เป็นบรรทัดเดียวต่อ Event ใช้ Token น้อยที่สุดแต่ข้อมูลบางส่วนหายไป)
นอกจากนี้ยังมีรูปแบบ \textbf{Summary} ซึ่งให้ LLM แปลง Logs เป็นคำอธิบายพฤติกรรมภาษาธรรมชาติก่อน โดยไม่ระบุชื่อหรือรหัสเทคนิค เพื่อใช้เป็น Input หรือ Query สำหรับค้นหาในฐานความรู้

% \section{ความสัมพันธ์เชิงกราฟของ Process}
% ความสัมพันธ์ระหว่าง Process สามารถแทนด้วยกราฟมีทิศทาง $G=(V,E)$
% โดยแต่ละโหนด $v \in V$ คือ Process (ระบุด้วย ProcessGuid)
% และเส้นเชื่อม $e \in E$ แทนความสัมพันธ์ เช่น Parent สร้าง Child (Sysmon EID 1)
% หรือ Process หนึ่งเข้าถึงอีก Process (EID 10)
% การท่องกราฟ (Graph Traversal) แบบ Breadth-First Search ที่จำกัดจำนวนขั้น (Hop)
% ทำให้สามารถรวบรวม Events ที่เชื่อมโยงกับ Alert ได้โดยไม่ลามไปทั้งระบบ

\section{ความสัมพันธ์เชิงกราฟของ Process}
ความสัมพันธ์ระหว่าง Process สามารถแทนด้วยกราฟมีทิศทาง $G=(V,E)$
โดยแต่ละโหนด $v \in V$ คือ Process (ระบุด้วย ProcessGuid)
และเส้นเชื่อม $e \in E$ แทนความสัมพันธ์เชิงสาเหตุ ได้แก่
Parent สร้าง Child (Sysmon EID 1) และ Process หนึ่งสร้างไฟล์ที่ถูกรันเป็นอีก Process (EID 11 ตามด้วย EID 1)
ส่วน Process Access (EID 10) ใช้เป็นหลักฐานประกอบของโหนดที่เกี่ยวข้อง แต่ไม่ใช้เป็นเส้นทางในการท่องกราฟ
การท่องกราฟโดยหาเส้นทางสั้นที่สุด (Shortest Path) จากจุดที่ Alert ชี้และจำกัดระยะทางสูงสุด
ทำให้สามารถรวบรวม Events ที่เชื่อมโยงกับ Alert ได้โดยไม่ลามไปทั้งระบบ

\section{ตัวชี้วัดที่ใช้ประเมินผล}
\begin{itemize}
  \item \textbf{Accuracy} -- สัดส่วนกรณีที่ระบบระบุเทคนิคตรงกับ Ground Truth
        \begin{equation}
        \mathrm{Accuracy} = \frac{\text{จำนวนกรณีที่ตอบถูก}}{\text{จำนวนกรณีทั้งหมด}}
        \end{equation}
  \item \textbf{Precision, Recall และ Macro F1} -- วัดความแม่นและความครอบคลุมรายเทคนิค
        แล้วเฉลี่ยทุกเทคนิคอย่างเท่ากัน
        \begin{equation}
        P_t = \frac{TP_t}{TP_t+FP_t},\quad
        R_t = \frac{TP_t}{TP_t+FN_t},\quad
        \text{Macro-}F_1 = \frac{1}{|T|}\sum_{t\in T}\frac{2P_tR_t}{P_t+R_t}
        \end{equation}
  \item \textbf{Latency} -- เวลาตั้งแต่รับ Logs จนได้คำตอบ (End-to-end)
  \item \textbf{Token Cost} -- จำนวน Token ที่ส่งให้ LLM เปรียบเทียบก่อนและหลังผ่าน Pipeline
\end{itemize}

\section{\ifenglish%
\ifcpe CPE \else ISNE \fi knowledge used, applied, or integrated in this project
\else%
ความรู้ตามหลักสูตรซึ่งถูกนำมาใช้หรือบูรณาการในโครงงาน
\fi
}
\begin{itemize}
  \item \textbf{ความมั่นคงปลอดภัยของระบบคอมพิวเตอร์และเครือข่าย} -- ความเข้าใจพฤติกรรมการโจมตี
        Logs ด้านความปลอดภัย และกระบวนการสืบสวนเหตุการณ์
  \item \textbf{ระบบปฏิบัติการ} -- กลไกการสร้างและความสัมพันธ์ระหว่าง Process บน Windows และ Linux
  \item \textbf{โครงสร้างข้อมูลและอัลกอริทึม} -- การสร้างกราฟ การท่องกราฟด้วยการหาเส้นทางสั้นที่สุด (Dijkstra)
        และการจัดลำดับด้วยคะแนนภายใต้งบประมาณ Token
  \item \textbf{ฐานข้อมูลและการจัดการข้อมูล} -- การ Parse, Normalize และจัดเก็บข้อมูลหลายรูปแบบ
  \item \textbf{วิศวกรรมซอฟต์แวร์} -- การออกแบบสถาปัตยกรรมระบบ การแบ่งโมดูล และการพัฒนา Web Application
  \item \textbf{ปัญญาประดิษฐ์/การเรียนรู้ของเครื่อง} -- การประเมินผลด้วย Accuracy, Precision, Recall, F1
        และการแยกข้อมูลสำหรับพัฒนากับข้อมูลสำหรับทดสอบ
\end{itemize}

\section{\ifenglish%
Extracurricular knowledge used, applied, or integrated in this project
\else%
ความรู้นอกหลักสูตรซึ่งถูกนำมาใช้หรือบูรณาการในโครงงาน
\fi
}
\begin{itemize}
  \item การติดตั้งและใช้งาน Splunk Enterprise รวมถึง SPL และ REST API
  \item รูปแบบ Logs ของ Windows Sysmon และ Linux Auditd และมาตรฐาน CIM
  \item กรอบความรู้ MITRE ATT\&CK และรูปแบบข้อมูล STIX
  \item การใช้งาน LLM ภายในเครื่องผ่าน Ollama, Prompt Engineering และ Structured Output
  \item การสร้างระบบ RAG ด้วย Embedding Model และ Vector Search
\end{itemize}
